\documentclass[acp,manuscript,noline]{copernicus}

\usepackage{xspace}

\begin{document}

\title{Supplementary material for ``Integrating generative AI with physics-based simulations for probabilistic modelling of volcanic clouds''}

\Author[1][lmingari@csic.es]{Leonardo}{Mingari}
\Author[1]{Arnau}{Folch}
\Author[1]{Heribert}{Pascual}
\Author[2]{Manuel}{Titos}

\affil[1]{Geosciences Barcelona (GEO3BCN-CSIC), Barcelona, Spain}
\affil[2]{University of Granada, Granada, Spain}

\runningtitle{Supplement: Generative AI for volcanic cloud modelling}
\runningauthor{L. Mingari}

\received{}
\pubdiscuss{egusphere-2025-6097}
\revised{}
\accepted{}
\published{}

\firstpage{1}

\maketitle
\tableofcontents

\section{Contents of the supplementary material}

This supplementary material accompanies the manuscript entitled ``Integrating
generative AI with physics-based simulations for probabilistic modelling of
volcanic clouds'' submitted to \textit{Atmospheric Chemistry and Physics}
(ACP), Copernicus manuscript reference egusphere-2025-6097. It provides
additional information on the ensemble configuration used in the numerical
experiments and lists the ancillary files and repositories required to
reproduce the simulations and figures.

\section{Ensemble configuration}

\subsection{Sampling technique}
Each ensemble was simulated by perturbing the Eruption Source Parameters (ESPs)
and wind components, widely recognised as some of the most sensitive
parameters/inputs for volcanic ash transport models \citep{devenish2012,poulidis2018}.
%
To efficiently sample the parameter space, we used Latin Hypercube Sampling
\citep[LHS,][]{mckay1979}, a space-filling ensemble technique that requires
fewer samples than purely random Monte Carlo sampling to achieve comparable
coverage. LHS was used to generate combinations of source-term and
meteorological perturbations for the FALL3D simulations.

\subsection{Synthetic experiment}

Table~\ref{tab:ensemble-synthetic} lists the reference value and sampling range
for each perturbed parameter. The top height of the eruption column was sampled
around the central value 9~\unit{km} above the vent level with a perturbation
range of $30\%$ and was assumed to be constant during the simulations for each
ensemble member. The vertical distribution of the source term was also modified
by perturbing the Suzuki parameter $A_s$ controlling the maximum emission
height \citep[see][for more information]{folch2020}. In addition, the
horizontal wind components were perturbed considering a sampling range of
$25\%$. The corresponding FALL3D model configuration is provided in the
supplementary file \texttt{garachico2025.inp}.

\begin{table}[h]
  \caption{
    Ensemble configuration for the synthetic experiment. The perturbed model
    parameters are: eruption column height ($H$), mass eruptive rate ($MER$),
    Suzuki parameter $A_s$, and the wind horizontal components.
  }
    \label{tab:ensemble-synthetic}
    \begin{tabular}{lccc}
    \tophline
    Parameter        & Reference value          & Distribution & Sampling range  \\ \middlehline
    $H$              & 9~km avl$^\dag$          & Uniform      & $\pm$30\%       \\
    $MER$            & $\sim 1.6\times 10^6$~kg~s$^{-1}$$^\ddag$ & Uniform      & $\pm$25\%       \\
    $A_s$            & 4                        & Uniform      & 1               \\
    Winds            & GFS$^\S$                 & Uniform      & $\pm$25\%       \\
    \bottomhline
    \end{tabular}
    \belowtable{%
        $^\dag$ Height given in km above the vent level; \\
        $^\ddag$ Estimated from $H$ according to~\citet{woodhouse2016}; \\
        $^\S$ Global Forecast System (GFS) weather forecast from NCEP.
    }
\end{table}

\subsection{Etna experiment}

For the Etna 24 December 2018 case, each ensemble was simulated by perturbing
the eruptive column-top height and the horizontal wind components by $20\%$. As
in the synthetic experiment, the perturbations were generated using LHS to
obtain an efficient sampling. A time-dependent eruptive column height was
defined in the simulations, as shown in Fig.~\ref{fig:etna}. The corresponding
FALL3D model configuration is provided in the supplementary file
\texttt{etna2018.inp}.

\begin{figure}[htbp]
  \centering
  \includegraphics[width=0.9\textwidth]{f01}
  \caption{
  Eruptive phases defined for the Etna experiment. 
  The simulations assumed a time-dependent eruptive column height
  (in meters above vent level).
  }
  \label{fig:etna}
\end{figure}

\section{Ensemble mean of the test dataset}
The test dataset used in the synthetic experiment consists of 2048 physics-based numerical simulations and is used to assess the performance of the VAE. Figure~\ref{fig:test} shows the ensemble mean of the test dataset, which is used as a reference for the computation of several metrics. The figure also indicates the locations of the points used to construct the histograms.

\begin{figure}[htbp]
  \centering
  \includegraphics[width=0.6\textwidth]{f02}
  \caption{
    Ensemble mean of the test dataset for the synthetic experiment, consisting of 2048 physics-based numerical simulations generated by the FALL3D dispersal model.
  }
  \label{fig:test}
\end{figure}

\section{Additional files and external repositories}

\subsection{Model configuration files}

The FALL3D model input files distributed with this supplementary material are:
\begin{itemize}
  \item \texttt{garachico2025.inp}: model configuration file for the synthetic experiment;
  \item \texttt{etna2018.inp}: model configuration file for the Etna 24 December 2018 experiment.
\end{itemize}

\subsection{Code and archived repositories}
The scripts used to post-process the ensemble simulations, VAE training,
ensemble sampling, and generate the figures are available from archived
repositories:
\begin{itemize}
  \item Synthetic experiment code Zenodo DOI \texttt{10.5281/zenodo.22759580} \citep{mingari2026supplementary-synthetic}; repository: \\
    \url{https://github.com/lmingari/egusphere-2025-6097-synthetic}
  \item Etna experiment code Zenodo DOI \texttt{10.5281/zenodo.22759573} \citep{mingari2026supplementary-etna}; repository: \\
    \url{https://github.com/lmingari/egusphere-2025-6097-etna}
\end{itemize}

\bibliographystyle{copernicus}
\bibliography{../references}

\end{document}
